\part{Using Git on your own}

\section{Using Git on your own}

\begin{frame}
\frametitle{Using Git on your own}
\begin{itemize}
    \item Why would you want to use Git alone?
    \item You obtain benefits of version control, e.g.:
        \begin{itemize}
            \item Have a safety net and access to a time machine.
            \item Get a log of changes over time and why.
            \item In the future, you can return to previous work and
                understand what your thoughts and ideas were.
        \end{itemize}
    \item You can use Git on your own computer, hence:
        \begin{itemize}
            \item setup and management is quick and simple,
            \item there are no servers to administer,
            \item no cloud service is necessary.
        \end{itemize}
    \item There are even more benefits when working with others; we'll
        discuss this in more detail later in the course.
\end{itemize}
\end{frame}

\section{Example project: a kiwi café}

\subsection{Getting started}

\begin{frame}{Example project: kiwi café}
Let's discover how Git can be useful with the help of a concrete example.

    \textbf{Introducing:}
    \begin{quote}
        Kiwi Cakes and Coffee\\
        Great coffee, awesome cake.
    \end{quote}

    \textbf{Scenario:}\\
    You're the proud owner of a new café with a focus on cake recipes found
    in New Zealand.
\end{frame}

\begin{frame}{Initial thoughts}
Where to start?
    \begin{itemize}
        \item The café is a big project.
        \item You'll want to collect the cake recipes
            somewhere sensible on your computer.  You and your staff
            can then find them easily.
        \item You'll probably want to store other information about your
            café close to the cake recipes.
        \item Thus, it makes sense to create a directory for the café and a
            subdirectory for the cake recipes.
    \end{itemize}
\end{frame}

\begin{frame}[fragile]{Project structure}
Create the project directories by running the following commands:
\begin{minted}{shell}
    mkdir kiwi-cakes-and-coffee
    cd kiwi-cakes-and-coffee
    mkdir recipes
\end{minted}

You now have this directory structure:
\begin{minted}{console}
    kiwi-cakes-and-coffee/
    └── recipes
\end{minted}

    Great!  A well-organised directory layout is a good first step.  Now we
    need to add some recipes.
\end{frame}

\subsection{A first recipe}

\begin{frame}[fragile]{A first recipe}
    Let's start with a classic that works anywhere: chocolate cake.

    \begin{itemize}
        \item Enter the \code{recipes} directory
        \begin{minted}{shell}
            cd recipes
        \end{minted}
        \item Open a new file in your text editor called \code{chocolate-cake.md}.
    \end{itemize}
\end{frame}

\begin{frame}[fragile]{A first recipe}
    Fill the file with this text
    (\href{https://raw.githubusercontent.com/paultcochrane/git_course/refs/heads/master/git/examples/edmonds-chocolate-cake.md}{\beamerbutton{link}}):
    \inputminted[fontsize=\tiny]{markdown}{examples/edmonds-chocolate-cake.md}
    {\footnotesize \emph{Edmonds Cookery Book}, 51st De Luxe Edition, 2003.}
\end{frame}

\begin{frame}{Why use Git?}
    \begin{itemize}
        \item We're likely to change recipes in the future.
        \item We'll also want to add and remove recipes.
        \item We'll need a system to store the files and manage the changes: Git.
        \item Thus, we need to create a Git repository in
            \code{kiwi-cakes-and-coffee} directory.
    \end{itemize}
\end{frame}

\subsection{Side note: concerning repositories}

\begin{frame}{What is a repository?}
    But first, what \emph{is} this repository thing?

    \begin{description}
        \item[repository] /ripózzitəri, -tri/, \emph{n.} a place where
            things are stored or may be found [\ldots].
            \emph{The Oxford Dictionary and Thesaurus, 1997.}
    \end{description}

    In short, it's a place to store stuff.

    In practice, it's a directory with some extra metadata.

    What does Git store in a repository?
    \begin{itemize}
        \item file content
        \item changes to file content
        \item references to file content
        \item references to other repositories
    \end{itemize}
\end{frame}

\subsection{Creating the Git repository}

\begin{frame}[fragile]{Creating the repository}
    Now that we know \emph{why} we want to store our recipes in a Git
    repository, let's create it.

    To create a Git repository, enter the project's base directory and run
    \code{git init}.

    \begin{minted}{shell}
        cd kiwi-cakes-and-coffee
        git init
    \end{minted}

    You should see this output:
    \begin{minted}[fontsize=\tiny]{output}
	hint: Using 'master' as the name for the initial branch. This default branch name
	hint: is subject to change. To configure the initial branch name to use in all
	hint: of your new repositories, which will suppress this warning, call:
	hint:
	hint:   git config --global init.defaultBranch <name>
	hint:
	hint: Names commonly chosen instead of 'master' are 'main', 'trunk' and
	hint: 'development'. The just-created branch can be renamed via this command:
	hint:
	hint:   git branch -m <name>
	Initialized empty Git repository in /home/vagrant/kiwi-cakes-and-coffee/.git/
    \end{minted}
\end{frame}

\begin{frame}[fragile]{\texttt{master} or \texttt{main} or \texttt{trunk} or \ldots?}
    That was a lot of output!
    \begin{itemize}
        \item Probably more than you expected.
        \item Important bit came at the end:
            \begin{minted}[fontsize=\tiny]{output}
	        Initialized empty Git repository in /home/vagrant/kiwi-cakes-and-coffee/.git/
            \end{minted}
        \item We've created our Git repository. \emoji{tada}
    \end{itemize}
\end{frame}

\begin{frame}[fragile]{\texttt{master} or \texttt{main} or \texttt{trunk} or \ldots?}
    What's with the hint text?
    \begin{itemize}
        \item Historically, the ``main'' branch was called \code{master}
        \item It's still the default name, but this may change in the
            future.
        \item Nowadays \code{main} is more common.
    \end{itemize}

    Change the branch name to \code{main}:
    \begin{minted}{shell}
        git branch -m main
    \end{minted}
\end{frame}

\subsection{Checking our surroundings}

\begin{frame}[fragile]{What's changed?}
    Initialising a Git repository creates a ``hidden'' directory called
    \code{.git}.

    Looking at the files and directories:
    \begin{minted}{console}
        kiwi-cakes-and-coffee/
        ├── .git
        │   <metadata and subdirs>
        └── recipes
            └── chocolate-cake.md
    \end{minted}

    The \code{.git} directory contains metadata and files about the
    repository.

    The files and directories outside the \code{.git} directory are called
    the \emph{working copy}.

    The repository itself is contained in the \code{.git} directory.
\end{frame}

\begin{frame}[fragile]{What state is the repository in?}
    \begin{itemize}
        \item A repository changes over time.
        \item Files are added, changed, deleted, moved, renamed.
        \item Use \code{git status} to find the current repository state:
    \end{itemize}
    \begin{minted}{shell}
        git status
    \end{minted}
    \begin{minted}[fontsize=\tiny]{output}
        On branch main

        No commits yet

        Untracked files:
          (use "git add <file>..." to include in what will be committed)
                recipes/

        nothing added to commit but untracked files present (use "git add" to track)
    \end{minted}
\end{frame}

\begin{frame}[fragile]{Lots of useful information}
    Git provides lots of information:
    \begin{minted}[fontsize=\tiny]{output}
        On branch main

        No commits yet

        Untracked files:
          (use "git add <file>..." to include in what will be committed)
                recipes/

        nothing added to commit but untracked files present (use "git add" to track)
    \end{minted}

    \begin{itemize}
        \item Where we are: \code{On branch main}
        \item What's already in the repository: \code{No commits yet}
        \item Files in the working copy: \code{Untracked files}
        \item How to add untracked files to the repository: \code{use "git
            add" to track}
        \item \emoji{rocket} Tip: read such output carefully.  All the help
            you need is often already provided in Git's messages.
    \end{itemize}
\end{frame}

\subsection{Adding our first recipe}

\begin{frame}[fragile]{Adding our first recipe}
    Remember the \code{git status} output from earlier?
    \begin{minted}[fontsize=\tiny]{output}
        Untracked files:
          (use "git add <file>..." to include in what will be committed)
                recipes/

        nothing added to commit but untracked files present (use "git add" to track)
    \end{minted}

    We want to track the chocolate cake recipe, hence we follow the advice
    in the status output and use \code{git add} to track it:
    \begin{minted}{shell}
        git add recipes/chocolate-cake.md
    \end{minted}
    This command produces no output; don't be surprised to see nothing
    appear on the console.
\end{frame}

\begin{frame}[fragile]{Advice when adding files}
    Note:
    \begin{itemize}
        \item We could have added the \code{recipes/} directory.
        \item \emph{In this case}, the effect would have been the same.
        \item In general: avoid adding/committing entire directories.
        \item Focus is your friend: which files do you really need?
        \item Be specific when adding files or file changes.
    \end{itemize}
\end{frame}

\begin{frame}[fragile]{Repository state after adding a file}
    What's the repository status now?

    How could we find this out?

    \only<2->{Use the \code{git status} command.  Try it out!}

    \begin{onlyenv}<3>
        \begin{minted}{shell}
            git status
        \end{minted}
        \begin{minted}{output}
            On branch main

            No commits yet

            Changes to be committed:
              (use "git rm --cached <file>..." to unstage)
                    new file:   recipes/chocolate-cake.md
        \end{minted}
    \end{onlyenv}

    \only<3>{We now have a new state: \code{Changes to be committed}}
\end{frame}

\subsection{Side note: committing and commits}

\begin{frame}[fragile]{Becoming committed}
    What's this commit concept?
    \begin{description}
        \item[commit] /kəmít/, \emph{v.tr.} entrust for safe
            keeping.\\
            \emph{The Oxford Dictionary and Thesaurus, 1997.}
    \end{description}

    For instance: \emph{to commit something to memory}.

    ``Commit'' has a connotation of more permanence than ``save''; like
    putting something into long-term storage.

    When we commit changes to a Git repository, we're putting them someplace
    safe, so that we can retrieve them later.

    A collection of changes committed at once is called \emph{a commit}.
\end{frame}

\subsection{Committing our first recipe}

\begin{frame}[fragile]{Committing our first recipe}
    We've only told Git to track our recipe by \emph{adding} it.

    To save this change permanently, we commit it with \code{git commit},
    e.g.:

    \begin{minted}{shell}
        git commit recipes/chocolate-cake.md
    \end{minted}

    This command will open a text editor with this content:
    \begin{minted}[fontsize=\tiny]{output}

        # Please enter the commit message for your changes. Lines starting
        # with '#' will be ignored, and an empty message aborts the commit.
        #
        # On branch main
        #
        # Initial commit
        #
        # Changes to be committed:
        #       new file:   recipes/chocolate-cake.md
        #
    \end{minted}

    \emoji{rocket} Tip: Again, Git provides lots of information.  Take the
    time to read and understand it.
\end{frame}

\subsection{Side note: editors for commit messages}

\begin{frame}[fragile]{Commit messages and text editors}
    Many command-line Git environments use Vim as the text editor for commit
    messages.  Some use a different editor, e.g.: Nano.

    You will need to know how to enter text into your given editor and how
    to save that text.

    This differs from environment to environment and editor to editor and is
    a common point of pain for newcomers to Git.

    Note: it is \emph{much} better long-term to use an editor than the
    \code{-m} command-line shortcut.

    \emoji{rocket} Tip: Take the time to choose an appropriate
    editor and/or learn the relevant commands for the editor you have.
\end{frame}

\begin{frame}[fragile]{Entering a commit message into Vim}
    To enter a commit message in Vim, follow these steps:
    \begin{itemize}
        \item To enter text, press \code{i} to change to \emph{input} mode.
        \item Enter the commit message text.
        \item To save the text, first press \code{Esc}, to change to
            \emph{command} mode.
        \item Then type \code{:wq} to \emph{write} the text to file and
            \emph{quit} the editor.
    \end{itemize}
\end{frame}

\begin{frame}[fragile]{Entering a commit message into Nano}
    To enter a commit message in Nano, follow these steps
    \begin{itemize}
        \item Enter the commit message text.
        \item To save the text, press \code{Ctrl} together with \code{x} to
            exit the editor.
        \item You will be asked if you wish to save the modified buffer.
        \item Press \code{y} to confirm saving.
        \item You will be shown the file name to write; this does not need
            to be changed.
        \item Press \code{Enter} to save the text and finish exiting the
            editor.
    \end{itemize}

    To use Nano instead of Vim in Git-Bash on Windows, run this command:
    \begin{minted}{shell}
        export EDITOR=nano
    \end{minted}
\end{frame}

\begin{frame}{Creating a new repository}
    Before we can add files and their content to a repository, we have to
    create/initialise one.

    Conceptual steps to create a Git repository:
    \begin{itemize}
        \item create a directory,
        \item change into that directory,
        \item initialise it!
    \end{itemize}

    The exact details depend upon the operating system.  We'll focus on the
    usage in a Bash-like shell.
\end{frame}

\begin{frame}[fragile]{Creating a new repository}
    Steps to create a Git repository (Bash-like shell):\footnote{For
    \code{cmd.exe} and Powershell the details are slightly different.}
    \begin{minted}{text}
        mkdir think-of-a-good-example-dir-name
        cd think-of-a-good-example-dir-name
        git init
    \end{minted}
    That's it!
\end{frame}

\begin{minted}{shell}
\end{minted}

\begin{frame}[fragile]{Exercise: create a new repository}
    Now it's your turn!
    \begin{itemize}
        \item Create a directory and change into it
    \end{itemize}
    \begin{minted}{shell}
    mkdir blupp
    cd blupp
    \end{minted}
    \begin{itemize}
        \item Initialise the repository
    \end{itemize}
        \begin{minted}{shell}
        git init
        \end{minted}
    \begin{itemize}
        \item What's the output of the \code{git init} command?  How does this
            help us?
        \item A directory was created (see command line output), what's in there?
        \item Use e.g.\ \code{ls -la} to find the hidden directory.
        \item What files are there in the hidden directory?
        \item Run \code{git status} for more info from git.  What's this
            output it telling us?
        \item Do you have any questions?
    \end{itemize}
\end{frame}

\begin{frame}[fragile]
    \inputminted{bash}{create-example-repo.txt}
    \inputminted{bash}{dot-git-dir-contents.txt}
\end{frame}

\begin{frame}[fragile]
\frametitle{What does git init do?}
\begin{itemize}
    \item Creates a \code{.git} directory
    \item Creates various configuration files, e.g. \code{.git/config}
    \item Sets up the repository; the repository is inside the \code{.git}
        directory
\end{itemize}
\end{frame}

% \begin{frame}[fragile]
% \begin{minted}{shell}
% $ git status
% On branch master
%
% Initial commit
%
% nothing to commit (create/copy files and use "git add" to track)
% \end{minted}
% \end{frame}
% %stopzone


\begin{frame}[fragile]
\frametitle{Adding files or changes to the repository}
\begin{itemize}
    \item To add files (or changes to files) to the repository's
        staging area, we use the \code{git add} command.
    \begin{itemize}
        \item the \emph{staging area} is where changes to the repository
            are kept before they are committed.
    \end{itemize}
    \item Git now knows to \emph{track} the file for future changes.
    \item \emph{Staging}, \emph{tracking}, and \emph{adding} are all
        roughly equivalent terms for the same process.
\end{itemize}

Let's add a file in tiny steps to see the process in detail:
    \begin{minted}{shell}
    touch moo     # create a file
    git status    # check repo status (file is *untracked*)
    git add moo   # add file to the staging area
    git status    # file is now a change to be committed
    \end{minted}
\end{frame}


\begin{frame}[fragile]
\frametitle{Committing files or changes to the repository}
\begin{itemize}
    \item Committing files or changes saves the staged changes into the
        repository.
    \item A commit is a snapshot of the current repository state.
\end{itemize}
    \begin{minted}{shell}
    git status    # file is now a change to be committed
    git commit    # commit staged changes to repository
    # enter commit message
    # see brief commit information
    git status    # working directory now clean
    \end{minted}
\end{frame}


\begin{frame}
\frametitle{Add/commit cycle}
The add/commit cycle is like how a photographer creates a group photo.
\end{frame}


\begin{frame}
\frametitle{Starting a new repository from scratch}
\begin{itemize}
    \item mkdir something
    \item cd something
    \item git init  (also ls -la)
    \item git status
    \item create file
    \item git status (untracked)
    \item git add
    \item git status (staged)
    \item git commit
    \item git status (clean)
    \item master is default branch
    \item (normally don't have to run git status all the time)
\end{itemize}
\end{frame}

\begin{frame}
    \frametitle{Starting a new repository (conceptual steps)}
    \begin{itemize}
        \item Create directory for repository (mkdir)
        \item Initialise repository (git init)
        \item Create file(s) (edit file)
        \item Add file(s) to repository (git add)
        \item Commit file(s) (git commit)
    \end{itemize}
\end{frame}

\begin{frame}
\frametitle{Starting a new repository (conceptual steps)}
\begin{itemize}
    \item To create a repository, need to \ttt{init}ialise it
        \begin{itemize}
            \item \code{git init}
        \end{itemize}
    \item Repository is just a directory, with a \code{.git} directory
        \begin{itemize}
            \item \code{ls -la}
        \end{itemize}
    \item Things that go into a repository are just files
        \begin{itemize}
            \item \code{touch file}
        \end{itemize}
    \item To tell Git to track a file, one needs to \ttt{add} it
        \begin{itemize}
            \item \code{git add}
        \end{itemize}
    \item Tracked files and \emph{changes to tracked files} are put in
        the \emph{staging area}
        \begin{itemize}
            \item \code{git status} (see staged files)
        \end{itemize}
    \item To record files and \emph{changes to files}, one \ttt{commit}s
        them to the repository
        \begin{itemize}
            \item \code{git commit}
        \end{itemize}
\end{itemize}
\end{frame}

\begin{frame}
\frametitle{Exercise: Start a new repository}
\begin{itemize}
    \item Repeat the previous steps on your own computer
    \item Create a directory and change into it
    \item Initialise the repository
    \item Run git status to see the current repository state
    \item Create a file (untracked); see the repository state, what changed?
    \item Add the file to the repository (track it); see the repository
        state, what changed?
    \item Commit the file to the repository; see the repository state, what
        changed?
\end{itemize}
\end{frame}

\begin{frame}
\frametitle{Getting Help}
\begin{itemize}
    \item man pages
    \item git help
    \item books
    \begin{itemize}
        \item How Git Works; Julia Evans; \url{https://wizardzines.com/zines/git/}
        \item Pro Git; Scott Chacon
    \end{itemize}
    \item online references
    \begin{itemize}
        \item Git book: \url{https://git-scm.com/doc}
        \item Cheat sheet: \url{https://github.com/dasesu/git/blob/main/Git_quick_reference.pdf}
    \end{itemize}
\end{itemize}
\end{frame}

\begin{frame}
\frametitle{Exercise: git help}
\begin{itemize}
    \item Use \code{git help} on \code{init}, \code{add}, \code{commit} and
        \code{config} commands
    \item Open and browse on \url{https://git-scm.org}
\end{itemize}
\end{frame}

\begin{frame}[fragile]
\frametitle{Parts of a local Git repository}
\begin{columns}[T]
    \begin{column}{0.5\textwidth}
        \begin{itemize}
            \item Working directory
            \item Staging area
            \item Repository
            \item Very important for understanding operations
        \end{itemize}
    \end{column}

    \begin{column}{0.5\textwidth}
        \begin{center}
            \resizebox{!}{0.7\textheight}{%
                \input{figures/working_dir_staging_repo.tex}
            }
        \end{center}
    \end{column}
\end{columns}
\end{frame}

\begin{frame}
\frametitle{Snapshots and diffs}
\begin{itemize}
    \item Photograph analogy
    \item Commits build on one another
    \item Working directory: checked out version of local repository
    \item Staging area: what will be committed next
    \item Repository: storage of snapshots
    \item Contrast with other systems, which store diffs between changes
\end{itemize}
\end{frame}

\begin{frame}
\frametitle{Basic workflow}
\begin{itemize}
    \item creation (mkdir, init, edit)
    \item saving snapshots of work (add, commit, return to edit step)
\end{itemize}
\end{frame}

\begin{frame}
\frametitle{Workflow}

Preparation:
\begin{itemize}
    \item \code{git init} (only once)
    \item \code{git checkout}
    \item \code{git config}
\end{itemize}

Main steps:
\begin{itemize}
    \item \code{git add} (before committing a file)
    \item \code{git mv}
    \item \code{git status}
    \item \code{git diff}
    \item \code{git rm}
    \item \code{git log}
    \item \code{git show}
    \item \code{git commit}
\end{itemize}
\end{frame}

\begin{frame}
\frametitle{Specifying commits}
\begin{description}
    \item[\code{c82a22c39c}] SHA hash
    \item[\code{HEAD}] name of the commit
    \item[\code{HEAD\^}] parent of \code{HEAD}
    \item[\code{HEAD\^\^}] grandparent of \code{HEAD}
    \item[\code{HEAD~2}] also grandparent of \code{HEAD}
    \item[\code{HEAD~4}] great-great grandparent of \code{HEAD}
\end{description}

also works with branch names:
\begin{description}
    \item[\code{main\^}] parent of tip of “\code{main}” branch
    \item[\code{main~4}] great-great grandparent of tip of “\code{main}” branch
        experimental tip of “\code{experimental}” branch
\end{description}
\end{frame}

\begin{frame}
\frametitle{Detour: Commit advice}

when, what, how to commit

commit messages

commit a single, focussed idea/concept

small commits have small diffs and are easier for others to review

Recommendations for commits:

\begin{itemize}
    \item A commit should contain a single, self contained idea.
    \item The commit message should contain a short description
        of the idea or change being made in this commit.
    \item A one line subject line
    \item An optional text describing more details of the change. The
        why of the change is important. Remember: this is a
        communication exercise.
    \item Commits should be as small as possible (atomic commits).
\end{itemize}

    http://sethrobertson.github.io/GitBestPractices

How often should I commit?
\begin{itemize}
    \item as often as possible
    \item “release early, release often”
    \item as soon as the changes belonging to a given idea are
        finished
    \item don’t be afraid to commit!
\end{itemize}

\end{frame}

\begin{frame}
\frametitle{Commit advice (cont.)}

\begin{itemize}
    \item Each commit needs a “commit message” describing the
    change
    \item Commit messages document what was done in each
    change and why
    \item An important part of communication within a project:
    \begin{itemize}
        \item commit logs are messages to other people
        \item \ldots and are messages to yourself in the future!
    \end{itemize}
\end{itemize}

What should I write in a commit message?
\begin{itemize}
    \item a short description of the change
    \item the description can sometimes be longer than the actual
    change! The message describes the “why” of the change.
    \item the message should be able to explain later what
    happened in the change
    \item note: when problems occur, this is where one looks first for
        information
\end{itemize}

\end{frame}

\begin{frame}
\frametitle{Extending the sample project}
\begin{itemize}
    \item edit, add, commit
    \item seeing what we've done
    \begin{itemize}
        \item git log (--graph)
        \item git show
        \item git diff
    \end{itemize}
    \item deleting/renaming files
    \begin{itemize}
        \item git rm
        \item git mv
    \end{itemize}
\end{itemize}
\end{frame}

\begin{frame}
\frametitle{Explain SHAs}
\begin{itemize}
    \item in context of diff, show, log, etc.
\end{itemize}
\end{frame}

\begin{frame}
\frametitle{Exercise: extend the sample project further}
\begin{itemize}
    \item recommended steps \ldots
    \item practice git add/commit/rm/mv
    \item practice git log/show/diff
\end{itemize}
\end{frame}

\begin{frame}
\frametitle{Ignoring files}

To ignore files in Git one simply needs a \code{.gitignore} file with
entries for the files one wishes to ignore. It is possible to use
regular expressions and wildcards such as \code{*.o} in order to ignore
multiple files.

TODO: insert example of creating files to ignore and then creating a
\code{.gitignore} file.
\end{frame}

\begin{frame}[fragile]
\frametitle{Ignoring files (cont.)}

Because the \code{.gitignore} file can change (and it’s highly likely
that we’re interested in these changes) one also keeps this file under
version control. We therefore add the file to the repository and commit
it:

    \inputminted{bash}{add-gitignore.txt}
\end{frame}

\begin{frame}[fragile]
\frametitle{Showing branches}
\code{git show-branch}

\begin{itemize}
    \item shows the available branches
    \item shows with an asterisk \code{*} which is the current branch
\end{itemize}

    \inputminted{bash}{git-show-branch-output.txt}
\end{frame}

\begin{frame}[fragile]
\frametitle{Showing branches (cont.)}
\code{git branch} shows available branches and which branch is
current.

    \inputminted{bash}{git-branch-output.txt}

\code{git branch -v} shows the commit hash and the commit
message of the most recent commit.

TODO need more expansive examples here with multiple branches.

    \inputminted{bash}{git-branch-v-output.txt}
\end{frame}

\begin{frame}
\frametitle{Import existing project}
\begin{itemize}
    \item git init, git add, git commit
\end{itemize}
\end{frame}

\begin{frame}
\frametitle{Aliasing commands}
\begin{itemize}
    \item reduce typing, use alias
    \item handy aliases: st, di, ci, co
\end{itemize}
\end{frame}

\begin{frame}
\frametitle{Staging and unstaging files}
\begin{itemize}
    \item staging area in more depth
    \item unstage files one doesn't want to track
    \item git reset (just moves pointer/label)
\end{itemize}
\end{frame}

\begin{frame}
\frametitle{Exercise!}
\end{frame}
